%%%PAGE 126 rot=0 legibility=good summary=球壳内部引力为0(同心球图、$g'$ 式、g=4/3πρGr 正比、g-r 图像)；挖空球体(割补法，在O2处的m受剩余部分引力)；页面下半约60%空白(仅透出背面反字)
%%%BLOCK id=126-01 ch=05 sec=球壳内部引力·g-r图像 kind=knowledge alt=none cont=none y=0.07-0.43
% 页面左缘/上缘露出邻页边缘(“No”“变”等)，下半页透出背面淡淡的反字，均已忽略
\textbf{球壳内部引力为0}

%FIG p126_a 0.14 0.15 0.25 0.225
\notefig[0.2]{figures/p126_a.png}
% 图：同心圆(外圆为球壳，内圆半径 $r$)，一条细通道自外圆通到内圆，质点 $m$ 与距离 $d$，另标 $R$、$r$ 等

$\dfrac{G\rho\,\frac{4}{3}\,r^{3}m}{d^{2}}=mg'$% CHECK: 分子原稿未见 $\pi$

$g=\dfrac{4}{3}\pi\rho g\,\textcircled{$r$}\propto r$% CHECK: 原稿 ρ 后写作小写 g(应为万有引力常量 G)；r 外有一圆圈

%FIG p126_b 0.09 0.272 0.32 0.42
\notefig[0.3]{figures/p126_b.png}
% 图：g-r 图像，纵轴 $g$、横轴 $r$；0 至虚线处为直线上升(旁注“正比”)，虚线处为最大值(横轴下注“$R_{\text{星球}}$.”)，之后曲线下降(旁注“平方反比”，写在白色修正贴上)
%%%END
%%%BLOCK id=126-02 ch=05 sec=非球体万有引力计算 kind=method alt=none cont=none y=0.10-0.25
%FIG p126_c 0.495 0.11 0.67 0.222
\notefig[0.3]{figures/p126_c.png}
% 图：大圆(半径 $2R$，圆心 $O_1$，竖直半径线旁标“2R”)，内含与大圆内切的小圆空腔(圆心 $O_2$，半径 $R$)，$m$ 位于 $O_2$；红色虚线圆为红笔所添

挖空$O_2$\ 在$O_2$的$m$受剩下部分的引力

$F_{\text{剩}}=\dfrac{G\,\frac{1}{8}\,Mm}{R^{2}}+\begin{matrix}F_{\text{补}}\\ \|\\ 0\end{matrix}$% 原稿 F补 正下方竖写“‖”与“0”，即 F补=0；CHECK: F 的下标像“剩”(也像“31”)
%%%END
%%%PAGEEND 126
